\chapter{研究背景与问题定义}

\section{研究背景}
多测站空间定位在雷达探测、光电跟踪、应急搜救、无人系统协同感知等领域具有基础作用。工程系统常见的观测量包括方位角、俯仰角和斜距。相对于纯测角问题，增加测距观测可明显改善可观性；相对于纯测距问题，加入测角观测可以在站位几何受限时提供方向补充。因此，融合测角与测距的联合定位模型具有更高的鲁棒性和更好的估计精度。

然而，该问题在数学上是典型非线性参数估计问题，主要难点包括：
\begin{enumerate}
  \item 观测函数非线性，且角度函数存在周期性与奇异性；
  \item 观测噪声方差不一致，必须进行权重归一化；
  \item 初值不合理时迭代算法可能震荡或收敛到次优点；
  \item 结果不仅需要点估计，还需要可解释的精度评估指标。
\end{enumerate}

\section{问题定义}
设目标位置向量为
\[
\bm p=[x,y,z]^{\T}\in\R^3.
\]
第 $j$ 个测角站坐标为 $\hat{\bm q}^{(A)}_j$，给出观测 $(A_j^{obs},h_j^{obs})$；第 $k$ 个测距站坐标为 $\hat{\bm q}^{(\rho)}_k$，给出观测 $\rho_k^{obs}$。其中 $A$ 表示方位角，$h$ 表示俯仰角，$\rho$ 表示斜距。

本文研究目标为：在给定多组观测及其精度指标的条件下，估计目标位置 $\bm p$，并完成以下任务：
\begin{itemize}
  \item 构建统一加权最小二乘模型；
  \item 设计可稳定收敛的数值求解算法；
  \item 输出估计精度与置信区域；
  \item 对噪声变化、几何变化进行灵敏度分析。
\end{itemize}

\section{符号约定}
为便于后续推导，统一记号如下：
\begin{table}[H]
\centering
\caption{主要符号说明}
\begin{tabular}{ll}
\toprule
符号 & 含义 \\
\midrule
$\bm p=[x,y,z]^{\T}$ & 目标位置向量 \\
$\hat{\bm q}_j=[\hat{x}_j,\hat{y}_j,\hat{z}_j]^{\T}$ & 第 $j$ 个测站坐标 \\
$A_j(\bm p)$ & 目标相对测站的理论方位角 \\
$h_j(\bm p)$ & 目标相对测站的理论俯仰角 \\
$\rho_j(\bm p)$ & 理论斜距 \\
$F(\bm p)$ & 加权残差平方和目标函数 \\
$\nabla F(\bm p)$ & 目标函数梯度 \\
$\bm H(\bm p)$ & 目标函数海森矩阵 \\
$\bm A$ & 法化矩阵（近似信息矩阵） \\
$\bm A^{-1}$ & 协方差近似矩阵 \\
\bottomrule
\end{tabular}
\end{table}

\section{研究思路与技术路线}
本文采用“几何初值 + 非线性迭代 + 统计评估”的技术路线。首先利用三测站测距几何交会得到初始点，随后分别采用牛顿法和法方程迭代法求解最优估计，最后基于海森矩阵或法化矩阵构建协方差近似与置信椭球。

从程序结构看，模型函数、导数函数、求解器与结果输出相互解耦，便于验证与扩展。本文在此基础上做了系统的数学化整理，以保证论文内容与代码实现的一致性。

\section{本文结构安排}
全文结构如下：
\begin{enumerate}
  \item 第 2 章建立观测几何与残差模型；
  \item 第 3 章推导梯度与海森矩阵；
  \item 第 4 章介绍初值生成与几何定位子问题；
  \item 第 5 章详细说明两类迭代算法及收敛判据；
  \item 第 6 章给出基准实验及结果；
  \item 第 7 章给出不确定性评估与置信区域；
  \item 第 8 章讨论扩展模型与工程实现建议；
  \item 第 9 章总结全文并提出未来工作。
\end{enumerate}
